\begin{table}[ht]
\centering\small
\begin{tabular}{@{}lrrrl@{}}
\toprule
Domain & Variables & Degree & Operations & Comparison\\
\midrule
$\Z$ & 13 & 1,072 & 195 & previous 1,088 / 195\\
$\Z$ & 14 & 414 & 148 & previous 696 / 150\\
$\Z$ & 15 & 218 & 130 & previous 214 / 139\\
$\Z$ & \textbf{16} & \textbf{124} & \textbf{118} & previous 184 / 135\\
$\Z_{>0}$ & 12 & 70 & 98 & previous 70 / 100\\
$\Z_{>0}$ & 13 & 51 & 88 & previous 51 / 90\\
\bottomrule
\end{tabular}
\caption{New literal native-family polynomials. Previous entries are selected
relevant points in the inspected repository, not assertions of global optima.}
\end{table}
The 16-integer result lowers degree by $32.6\%$ and operations by $12.6\%$
against the degree-184, 135-operation parent. It also improves the separate
degree-192, 134-operation point in both measures. The 14-integer equation
has a lower degree than both the 696/150 root-fusion equation and the
652/230 fixed-prime equation. The 792/137 equation remains a lower-cost
alternative. The 15-integer entry is a degree--cost tradeoff.

\begin{table}[ht]
\centering\small
\begin{tabular}{@{}lrrrrr@{}}
\toprule
Direct quartic over $\N$ & Variables & Rows & Terms & Height & Degree\\
\midrule
Parent, affine elimination & 3,614 & 2,954 & 15,996 & 127,778,545 & 4\\
New, affine elimination & \textbf{2,473} & 2,089 & 20,580 & 127,808,289 & 4\\
Parent, retained affine rows & 6,183 & 5,523 & 17,430 & 1,993 & 4\\
New, retained affine rows & \textbf{4,860} & 4,476 & 14,719 & \textbf{1,837} & 4\\
\bottomrule
\end{tabular}
\caption{Fully expanded direct polynomials. Height is the maximum absolute
coefficient. Rows are quadratic residuals before summing their squares.}
\end{table}
The first new direct equation has $31.6\%$ fewer variables than its parent,
at the cost of more expanded monomials and a slightly larger coefficient
height. The retained-row variant reduces the variable, row, and monomial
counts and the coefficient height against its corresponding parent,
while retaining degree four. Its coefficients need at most eleven magnitude bits.

